\chapter[Sermon 98. St. John Is Commanded Not to Seal This Book, but to Publish It, Having Respect to No Man.]{St. John Is Commanded Not to Seal This Book, but to Publish It, Having Respect to No Man.}
\label{sermon:098}
\markright{Sermon 98. St. John Is Commanded Not to Seal This Book, but to ...}

And he said to me: Seal not the sayings of the prophecy of this book. For the time is at hand. He who does evil, let him do evil still; and he who is defiled, let him be defiled still; and he who is righteous, let him be more righteous; and he who is holy, let him be more holy.

\index{Antichrist}\index{Arethas of Caesarea}\index{Daniel (Book of)}7 The seventh point in this conclusion forbids John from sealing the book that was written. But the angel says, “Seal it not.” Certainly, letters and books are often sealed either for the sake of credibility and confirmation, or so that they are not openly read by all, but only by those to whom they are assigned. An angel says to Daniel in chapter 12, “You, Daniel, seal the words and the book until the last time.” He is commanded to shut the book, that is to say, to make an end, and to expect no further revelation. Finally, he is commanded to shut it for the ungodly, to whom this book will assuredly seem dark and closed. For it follows: “Many shall err, and knowledge shall be manifold.” For those not ruled by the certain and sure word of God have nothing certainly tried and known, but wander through manifold, sundry, and uncertain opinions, judgments, and traditions of men. For Daniel says that knowledge shall be variable: that is to say, there shall be innumerable opinions and sects concerning the religion and service of God; whereas there is but one only true opinion, doctrine, faith, or religion—the same that Daniel set forth in his book, which book also he sealed, that is to say, confirmed it as if with godly seals, as authentic and authorized, and worthy to be credited. However, at this present, John is not commanded in the same sense and meaning not to seal his book, which we know to be altogether authentic. But what the angel means by such a thing is to conceal or cover not, and hide not this book: so that it may be written, that it might be a public doctrine throughout the world, whereby all men might be instructed in the things revealed from heaven, that they be not, through the schemes and tyranny of Antichrist, withdrawn from the kingdom of Christ to the kingdom of Antichrist. For God would that all these things should be most common and manifestly known to all men. This sense has Arethas also explained, saying: “Seal them not,” that is, keep them not sealed to yourself, but publish them to all. The reason is added: for the time is at hand, wherein verily these things which I have said shall come to pass. Wherefore the faithful had need of warning, confirming, and comfort. Considering therefore that this book is set forth that it might admonish, strengthen, and comfort the faithful, the same ought not to be shut but wide open. For this is the good will of God, that this his word should be preached in his church to the profit of all faithful. Let them look, therefore, what they do, who would have this book not only shut up, but clean taken away; neither think it can be understood as obscure and full of dark speakings. But to God be praise and thanks giving, who has vouchsafed to provide for us faithfully and in time by this most profitable and most necessary book.

\index{Ezekiel (Book of)}The eighth point of this conclusion seems to address a certain objection. For someone might say: You want this book to be open and accessible to all people of all ranks, genders, and ages. But there will be some who will utterly despise it. Therefore, preaching it will be in vain, and we will urge these writings in vain, especially those who will mock it and interpret it as they please. But he seems to anticipate this, saying: Surely, there will be countless unrighteous people who will proceed unchecked in their wickedness and will increasingly exceed and surpass themselves. Yet, there will also be righteous people who, persevering in all righteousness, will grow in holy virtues and will excel in this as well. Therefore, do not hesitate to proclaim to all what I have commanded you in this book, without being concerned about the outcome. Leave that to me; fulfill the office of preaching. I will ensure that your faithful preaching is not in vain. And leave them alone, if you see that some will be entirely depraved and perish in their depravity, saying they despise all your faithful labor. For you have done your duty and are blameless. They perish through their own fault. Therefore, I do not want you or anyone else to be overly concerned when you see many rejecting the purity of God’s word and preferring to wallow in filth. And we read elsewhere,[?], that the Gospel is preached to many for their condemnation, and the fragrance of the gospel is sweet to some for salvation, but to others an intolerable sting leading to destruction. A similar passage is found in Ezekiel 2, where we read that the Lord said to the prophet: “Son of man, I am sending you to the people of Israel, a rebellious people who have rebelled against me, they and their fathers, until this day. They are a people of stubborn heart and a perverse mind. I am sending you to them, and you shall tell them: Thus says the Lord God, whether they hear or not, for it is a rebellious house, so that they may know that there has been a prophet among them. And you, son of man, do not fear them or be afraid of their words, for they are contentious and sharp like thorns, and you live among scorpions. But because of this, you must not be afraid of them; you must speak my words to them, whether they hear or not.”

\index{Aquinas, Thomas}However, we must take heed here that we do not understand that God commands the ungodly to proceed to be more ungodly, where the angel says, “Let him who is unjust be unjust still,” and so forth. For it seems to be a saying in a manner like this, as it is in the gospel: “Whatever you are doing, do it more speedily.” For he commands him to do that which he knew he would do. In the same way here also, that he knew the wicked would do, he says they shall do; neither does he will that their doings should trouble John and the faithful preacher, saying that there shall also be many good people who shall also apply themselves to righteousness. We are wont also to say with a much like phrase: if it will not be otherwise, we must be content. Not that we bid him who perishes to perish; but that so we reproach his madness to him, and signify that he perishes through his own fault, willingly and knowingly. Arethas says: “It is no exhortation,” but rather a rebuking of every one to whom he applies himself. And Thomas of Aquin says: “The sense is, let him who hurts, hurt still.” That is, he will hurt by doing other evils, so that the angel be understood to have said these things in prophesying, not in wishing, and so forth. And so the meaning would be: the wicked confirming the prophecy shall continue to be wicked, the godly again shall grow in the holy study of righteousness. This sense truly seems most plain of all. Neither do they differ much from these that are read in the twelfth chapter of Daniel by these words: “Go, Daniel,” says the angel, “and search not curiously over the instant of the last time; for the sayings are closed and sealed until the last time. Many shall be purified and made white and cast anew. But the wicked shall do wickedly, and all ungodly shall not understand, but the learned shall teach.” From these things, nothing at all swerves from the words of the Apostle, speaking and prophesying of the later times: “All who will live godly in Christ Jesus shall suffer persecution for righteousness. Nevertheless, evil men and deceivers grow worse and worse, while they both lead others into error and err themselves.” Therefore, saying that the later age of this world shall be such, let us, who are called to this function, proceed constantly to advance, set forth, and beat in, the very word of God and revelation of Jesus Christ unto all men, regarding nothing what the world and worldly men speak against it.

And he skillfully contrasts two kinds of men against two others, the unrighteous against the righteous, and the defiled against the holy. “Let him who does evil,” he says, “continue to do evil; let him who is unrighteous continue to be unrighteous; let him who harms by persecuting the godly continue to harm, and indeed, to intensify it.” Against this, he sets: “Let him who is righteous be even more righteous, let him advance further, and grow increasingly in all godliness, and go beyond himself in righteousness, both in faith and works.” For by the righteousness of faith we are justified; by the righteousness of works, we are declared to be righteous. And those who are righteous not only harm no one but also benefit and do good to all. Conversely, the unrighteous, who lack true faith, lack light, and therefore walk in darkness, doing the works of darkness, persecuting both the righteous, and righteousness, and troubling all. And that there should be such men in the later days, the Lord also prophesied in the twenty-fourth chapter of the Gospel according to Matthew.

Another kind of men are unclean, polluted, filthy, and vile, and so forth. He says, “Let the filthy be silent in their filth.” And the interpreters of the Greek language advise that [illegible] refers to that filth which we gather with the ends of our nails. He signified unclean persons in body and soul: idolaters, fornicators, gluttons, and the like. Against them he has placed the holy, pure, and clean—that is to say, purified by faith and diligently applying themselves to sanctification. Therefore, just as the filthy wallow themselves more and more in the mire, and array and defile themselves vilely, so the godly apply themselves more and more daily to cleanliness and holiness of life. The Lord Jesus justifies and sanctifies us forevermore.
